\documentclass[a4paper,10pt,brazil]{article}
\usepackage{provastyle}
\usepackage{menukeys}

\begin{document}
\makeexamtitlepage{UFOP}{BCC222: Programação Funcional}{José Romildo}{2026/1}{Primeira Prova}{09}
\begin{multicols}{2}
\questionTitle{1}{Vantagem de tipos sinônimos}
\begin{flushright}
  \footnotesize
  \textit{\textbf{4: Tipos de Dados}}\\
  \textit{c04-sinonimos-vantagem-01}\\

\end{flushright}
Qual é o principal benefício de usar um tipo sinônimo, como \haskellinline{type Idade = Int}?

\begin{enumerate}
  \item Permitir a definição de tipos recursivos.

  \item Melhorar o desempenho do programa, pois sinônimos são mais eficientes.

  \item É a única maneira de usar \haskellinline{Int} em assinaturas de função.

  \item Aumentar a legibilidade e a expressividade do código, tornando as assinaturas mais auto-documentadas.

  \item Criar um tipo completamente novo que é incompatível com \haskellinline{Int}, aumentando a segurança.


\end{enumerate}

\questionTitle{2}{Ausência de coerção implícita entre Int e Double}
\begin{flushright}
  \footnotesize
  \textit{\textbf{7: Polimorfismo}}\\
  \textit{cap07-conversao-int-double-soma}\\

\end{flushright}
Considere:

\begin{minted}{haskell}
n :: Int
n = 10

x :: Double
x = 2.5
\end{minted}

Qual expressão soma corretamente os valores como \haskellinline{Double}?

\begin{enumerate}
  \item \mintinline{haskell}{show n + x}
  \item \mintinline{haskell}{fromIntegral n + x}
  \item \mintinline{haskell}{toInteger x + n}
  \item \mintinline{haskell}{n + x}
  \item \mintinline{haskell}{read n + x}

\end{enumerate}

\questionTitle{3}{Modelagem de biblioteca}
\begin{flushright}
  \footnotesize
  \textit{\textbf{6: Estruturas de dados básicas}}\\
  \textit{c06-modelagem-01}\\

\end{flushright}
Considere as definições abaixo:
\begin{minted}{haskell}
type Titulo = String
type Autor  = String
type Ano    = Maybe Int
type Livro  = (Titulo, Autor, Ano)
\end{minted}
Qual definição representa corretamente uma biblioteca como coleção de livros?

\begin{enumerate}
  \item \haskellinline{type Biblioteca = Maybe Livro}

  \item \haskellinline{type Biblioteca = Livro}

  \item \haskellinline{type Biblioteca = (Livro, Livro)}

  \item \haskellinline{type Biblioteca = [Titulo]}

  \item \haskellinline{type Biblioteca = [Livro]}


\end{enumerate}

\questionTitle{4}{Declaração let no GHCi}
\begin{flushright}
  \footnotesize
  \textit{\textbf{3: Expressões e definições}}\\
  \textit{c03-ghci-definicoes-02}\\

\end{flushright}
No GHCi, qual alternativa representa a forma explícita de associar o nome \haskellinline|raio| ao valor \haskellinline|10| para uso em expressões posteriores?

\begin{enumerate}
  \item \haskellinline|let 10 = raio|.

  \item \haskellinline|let raio = 10|.

  \item \haskellinline|raio <- 10|.

  \item \haskellinline|raio := 10|.

  \item \haskellinline|def raio = 10|.


\end{enumerate}

\questionTitle{5}{Tipo de uma média com fromIntegral}
\begin{flushright}
  \footnotesize
  \textit{\textbf{7: Polimorfismo}}\\
  \textit{cap07-expressao-media-corrigida}\\

\end{flushright}
Qual é o tipo mais geral da expressão abaixo?

\begin{minted}{haskell}
fromIntegral (length [True, False]) / 2
\end{minted}

\begin{enumerate}
  \item \mintinline{haskell}{Fractional a => a}
  \item \mintinline{haskell}{[Bool] -> Int}
  \item \mintinline{haskell}{Int}
  \item \mintinline{haskell}{Integer}
  \item \mintinline{haskell}{Bool}

\end{enumerate}

\questionTitle{6}{Definição de tipo em Haskell}
\begin{flushright}
  \footnotesize
  \textit{\textbf{4: Tipos de Dados}}\\
  \textit{c04-conceito-definicao-01}\\

\end{flushright}
Em Haskell, um \textbf{tipo} pode ser melhor definido como:

\begin{enumerate}
  \item Uma coleção de valores que compartilham propriedades e operações comuns.

  \item Um comentário especial que descreve o propósito de uma função.

  \item Uma função que recebe argumentos e retorna um resultado.

  \item Uma palavra-chave reservada para declaração de constantes.

  \item Um nome dado a uma variável que armazena um valor numérico.


\end{enumerate}

\questionTitle{7}{Consulta com zip e lookup}
\begin{flushright}
  \footnotesize
  \textit{\textbf{6: Estruturas de dados básicas}}\\
  \textit{c06-lookup-associacao-01}\\

\end{flushright}
Considere a expressão abaixo.
\begin{minted}{haskell}
let chaves = [1,2,3,4]
    valores = ["a","b","c","d"]
    tabela = zip chaves valores
in lookup 3 tabela
\end{minted}
Qual é o resultado final?

\begin{enumerate}
  \item Ocorre um erro de tipo.

  \item \haskellinline{ Just "c" }

  \item \haskellinline{ Just "d" }

  \item \haskellinline{ Nothing }

  \item \haskellinline{ "c" }


\end{enumerate}

\questionTitle{8}{Status calculado localmente}
\begin{flushright}
  \footnotesize
  \textit{\textbf{5: Expressão condicional}}\\
  \textit{c05-where-avaliacao-01}\\

\end{flushright}
Analise a função:
\begin{haskellcode}
processa x
  | status == "par"   = x * 2
  | status == "impar" = x + 1
  where
    status = if even x then "par" else "impar"
\end{haskellcode}
Qual é o resultado de \haskellinline|processa 7|?

\begin{enumerate}
  \item \haskellinline|7|

  \item \haskellinline|14|

  \item \haskellinline|8|

  \item Ocorre um erro, pois \haskellinline|status| só estaria visível depois da cláusula \haskellinline|where|.

  \item Ocorre um erro de tipo na comparação com \haskellinline|"impar"|.


\end{enumerate}

\questionTitle{9}{Shadowing em expressões let aninhadas}
\begin{flushright}
  \footnotesize
  \textit{\textbf{3: Expressões e definições}}\\
  \textit{c03-let-escopo-02}\\

\end{flushright}
Considere a expressão:

\begin{haskellcode}
let x = 10
    y = x + 1
in let x = 3
   in x + y
\end{haskellcode}

Qual é o seu valor?

\begin{enumerate}
  \item \haskellinline|23|, pois os dois valores de \haskellinline|x| são somados a \haskellinline|y|.

  \item \haskellinline|6|, pois o \haskellinline|x| interno redefine também o valor de \haskellinline|y|.

  \item Ocorre erro de escopo, pois não é permitido reutilizar o nome \haskellinline|x| em um \haskellinline|let| aninhado.

  \item \haskellinline|14|, pois o \haskellinline|x| interno vale \haskellinline|3| no corpo interno, mas \haskellinline|y| foi definido a partir do \haskellinline|x| externo.

  \item \haskellinline|21|, pois o \haskellinline|x| externo continua valendo em todo o corpo interno.


\end{enumerate}

\questionTitle{10}{Significado de otherwise}
\begin{flushright}
  \footnotesize
  \textit{\textbf{5: Expressão condicional}}\\
  \textit{c05-otherwise-01}\\

\end{flushright}
O que \haskellinline|otherwise| representa em uma definição com guardas?

\begin{enumerate}
  \item Um nome predefinido equivalente a \haskellinline|True|, normalmente usado como caso padrão.

  \item Uma variável criada automaticamente com o valor da guarda anterior.

  \item Um nome reservado equivalente a \haskellinline|False|.

  \item Uma forma especial de indicar erro em tempo de execução.

  \item Uma palavra reservada que faz a avaliação voltar para a equação anterior.


\end{enumerate}

\questionTitle{11}{Sintaxe correta com guardas}
\begin{flushright}
  \footnotesize
  \textit{\textbf{5: Expressão condicional}}\\
  \textit{c05-guardas-sintaxe-semantica-01}\\

\end{flushright}
Qual das definições abaixo está escrita corretamente usando guardas em Haskell?

\begin{enumerate}
  \item \begin{haskellcode}
abs' n
  | n < 0 then -n
  | otherwise then n
\end{haskellcode}

  \item \begin{haskellcode}
abs' n
  | n < 0     = -n
  | otherwise = n
\end{haskellcode}

  \item \begin{haskellcode}
abs' n
  if n < 0 = -n
  else     = n
\end{haskellcode}

  \item \begin{haskellcode}
abs' n
  n < 0 -> -n
  otherwise -> n
\end{haskellcode}

  \item \begin{haskellcode}
abs' n =
  | n < 0     = -n
  | otherwise = n
\end{haskellcode}


\end{enumerate}

\questionTitle{12}{Boas práticas em assinaturas explícitas}
\begin{flushright}
  \footnotesize
  \textit{\textbf{7: Polimorfismo}}\\
  \textit{cap07-comparacao-assinaturas-explicitas}\\

\end{flushright}
Mesmo com inferência de tipos, por que é boa prática escrever assinaturas explícitas para funções de topo?
\begin{enumerate}
  \item Porque Haskell não consegue inferir tipos de funções definidas pelo programador.
  \item Porque a assinatura documenta a intenção, ajuda a localizar erros e pode impedir que uma função fique mais geral ou mais restrita do que o planejado.
  \item Porque toda função sem assinatura passa a ter tipo \haskellinline{String}.
  \item Porque assinaturas explícitas desativam a checagem de tipos, facilitando a compilação.
  \item Porque assinaturas explícitas tornam as funções mutáveis.

\end{enumerate}

\questionTitle{13}{Carregando módulos no GHCi}
\begin{flushright}
  \footnotesize
  \textit{\textbf{2: Ambiente de desenvolvimento}}\\
  \textit{c02-ghci-modulos-02}\\

\end{flushright}
Para utilizar funções avançadas de listas (como \haskellinline{sort}) diretamente no prompt do GHCi, o desenvolvedor precisa trazer o módulo \texttt{Data.List} para o escopo atual. Quais comandos são válidos para realizar esta operação?

\begin{enumerate}
  \item Não é necessário usar comandos; em Haskell, diferentemente de C ou Java, todos os módulos oficiais estão sempre injetados globalmente no Prelude.

  \item O comando \texttt{:browse Data.List}, que carrega e executa todas as funções contidas no módulo de forma simultânea.

  \item O comando \texttt{:install Data.List}, que solicita a função diretamente do repositório remoto Hackage em tempo real.

  \item Apenas o comando interno \texttt{:load Data.List}, pois o termo \haskellinline{import} é exclusivo para arquivos físicos \texttt{.hs}.

  \item A palavra-chave nativa da linguagem \haskellinline{import Data.List} ou o atalho específico do GHCi \texttt{:module +Data.List} (ou \texttt{:m +Data.List}).


\end{enumerate}

\questionTitle{14}{Restrição Num e ausência de coerção implícita}
\begin{flushright}
  \footnotesize
  \textit{\textbf{7: Polimorfismo}}\\
  \textit{cap07-adhoc-restricao-num}\\

\end{flushright}
Considere a assinatura:

\begin{minted}{haskell}
(+) :: Num a => a -> a -> a
\end{minted}

Qual interpretação é a mais precisa?

\begin{enumerate}
  \item O operador aceita livremente um \haskellinline{Int} e um \haskellinline{Double}, pois ambos pertencem à hierarquia numérica.
  \item A assinatura diz que \haskellinline{(+)} é paramétrica pura, sem qualquer restrição sobre \haskellinline{a}.
  \item Os dois operandos e o resultado têm o mesmo tipo \haskellinline{a}; esse tipo deve ser instância de \haskellinline{Num}.
  \item A restrição \haskellinline{Num a} significa que o resultado será sempre \haskellinline{Integer}.
  \item O compilador escolhe uma implementação por herança de classe de POO.

\end{enumerate}

\questionTitle{15}{Descarte de valores em do}
\begin{flushright}
  \footnotesize
  \textit{\textbf{8: Programas interativos}}\\
  \textit{c08-blocos-do-04}\\

\end{flushright}
No bloco abaixo, a primeira chamada a \haskellinline|getLine| está isolada em uma linha:

\begin{minted}{haskell}
main :: IO ()
main = do
  getLine
  nome <- getLine
  putStrLn nome
\end{minted}

Qual é o comportamento correto?

\begin{enumerate}
  \item A primeira chamada a \haskellinline|getLine| não é executada, pois seu resultado não é usado.

  \item A primeira ação \haskellinline|getLine| é executada e consome uma linha da entrada, mas seu resultado é descartado. A segunda leitura é associada a \haskellinline|nome| e depois impressa.

  \item O programa imprime as duas linhas lidas, uma em cada linha.

  \item As duas chamadas a \haskellinline|getLine| leem a mesma linha, pois Haskell reutiliza automaticamente ações de E/S.

  \item O programa não compila, pois toda ação em um bloco \haskellinline|do| deve usar \haskellinline|<-|.


\end{enumerate}

\questionTitle{16}{Comando :type no GHCi}
\begin{flushright}
  \footnotesize
  \textit{\textbf{4: Tipos de Dados}}\\
  \textit{c04-ghci-consulta-01}\\

\end{flushright}
Qual comando do GHCi exibe o tipo de uma expressão sem avaliá-la?

\begin{enumerate}
  \item \texttt{:info}

  \item \texttt{:load}

  \item \texttt{:show}

  \item \texttt{:type}

  \item \texttt{:kind}


\end{enumerate}

\questionTitle{17}{Assinatura de função com Bool e lista de Char}
\begin{flushright}
  \footnotesize
  \textit{\textbf{4: Tipos de Dados}}\\
  \textit{c04-funcao-assinatura-02}\\

\end{flushright}
Qual das seguintes assinaturas de tipo representa uma função que recebe um \haskellinline{Bool} e uma lista de \haskellinline{Char} e retorna um \haskellinline{Char}?

\begin{enumerate}
  \item \haskellinline{(Bool, [Char]) -> Char}

  \item \haskellinline{[Bool] -> [Char] -> Char}

  \item \haskellinline{(Bool, [Char], Char)}

  \item \haskellinline{Bool -> [Char] -> Char}

  \item \haskellinline{Char -> Bool -> [Char]}


\end{enumerate}

\questionTitle{18}{Ambiguidade de show [] e EDR}
\begin{flushright}
  \footnotesize
  \textit{\textbf{7: Polimorfismo}}\\
  \textit{cap07-defaulting-show-lista-vazia}\\

\end{flushright}
Considere a definição:

\begin{minted}{haskell}
w = show []
\end{minted}

Qual comparação entre compilação normal e GHCi com \haskellinline{ExtendedDefaultRules} ligada é correta?

\begin{enumerate}
  \item A expressão só compila se a lista vazia for convertida com \haskellinline{fromInteger}.
  \item A expressão sempre tem tipo \haskellinline{[String]}, pois \haskellinline{show} retorna listas.
  \item A expressão sempre é rejeitada, inclusive no GHCi com EDR, pois listas vazias não têm tipo.
  \item O Relatório escolhe \haskellinline{Integer} para o elemento da lista, pois \haskellinline{Integer} é o primeiro default.
  \item Pela regra do Relatório, a expressão é ambígua por não haver classe numérica; no GHCi com EDR, ela pode ser resolvida escolhendo \haskellinline{a ~ ()}.

\end{enumerate}

\questionTitle{19}{putStrLn e prompts}
\begin{flushright}
  \footnotesize
  \textit{\textbf{8: Programas interativos}}\\
  \textit{c08-buffering-04}\\

\end{flushright}
Por que \haskellinline|putStrLn "Digite seu nome:"| costuma evitar o problema de \emph{prompt} invisível que pode ocorrer com \haskellinline|putStr "Digite seu nome: "|?

\begin{enumerate}
  \item Porque \haskellinline|putStrLn| só funciona dentro de programas compilados, não no GHCi.

  \item Porque \haskellinline|putStrLn| executa automaticamente \haskellinline|getLine| depois de imprimir.

  \item Porque \haskellinline|putStrLn| não usa \emph{buffer}, enquanto \haskellinline|putStr| sempre usa.

  \item Porque \haskellinline|putStrLn| acrescenta uma quebra de linha; em muitos ambientes interativos, essa quebra de linha faz a saída ser descarregada. Porém, ela também coloca a entrada do usuário na linha seguinte.

  \item Porque \haskellinline|putStrLn| converte o texto para \haskellinline|IO String| antes de imprimir.


\end{enumerate}

\questionTitle{20}{Suficiência do tipo Int}
\begin{flushright}
  \footnotesize
  \textit{\textbf{4: Tipos de Dados}}\\
  \textit{c04-tipos-int-integer-02}\\

\end{flushright}
Para qual dos seguintes valores o tipo \haskellinline{Int} é geralmente suficiente, não necessitando de \haskellinline{Integer}?

\begin{enumerate}
  \item O número de grãos de areia na Terra.

  \item A idade de uma pessoa.

  \item A distância em milímetros entre duas galáxias.

  \item A população mundial.

  \item O resultado de \haskellinline{2 ^ 200}.


\end{enumerate}

\questionTitle{21}{Avaliação direta de if}
\begin{flushright}
  \footnotesize
  \textit{\textbf{5: Expressão condicional}}\\
  \textit{c05-if-avaliacao-01}\\

\end{flushright}
Qual é o valor da expressão abaixo?
\begin{haskellcode}
if 5 > 10 then 1 + 1 else 2 + 2
\end{haskellcode}

\begin{enumerate}
  \item \haskellinline|"4"|

  \item \haskellinline|4|

  \item \haskellinline|2|

  \item Ocorre um erro de sintaxe.

  \item \haskellinline|False|


\end{enumerate}

\questionTitle{22}{Decisão em programa interativo}
\begin{flushright}
  \footnotesize
  \textit{\textbf{8: Programas interativos}}\\
  \textit{c08-programas-completos-04}\\

\end{flushright}
Considere:

\begin{minted}{haskell}
main :: IO ()
main = do
  nota <- readLn :: IO Double
  if nota >= 6.0
    then putStrLn "Aprovado"
    else putStrLn "Reprovado"
\end{minted}

Se o usuário digitar \haskellinline|5.5|, o que será impresso?

\begin{enumerate}
  \item \haskellinline|5.5|

  \item \haskellinline|False|

  \item \haskellinline|Aprovado|

  \item \haskellinline|Reprovado|

  \item O programa não compila, pois \haskellinline|if| não pode aparecer dentro de \haskellinline|do|.


\end{enumerate}

\questionTitle{23}{Significado do operador ::}
\begin{flushright}
  \footnotesize
  \textit{\textbf{4: Tipos de Dados}}\\
  \textit{c04-assinatura-operador-01}\\

\end{flushright}
Na declaração \haskellinline{minhaFuncao :: Int -> Bool}, o que o operador \haskellinline{::} significa?

\begin{enumerate}
  \item \emph{é um subtipo de}

  \item \emph{tem o tipo de}

  \item \emph{pertence à classe}

  \item \emph{é definido como}

  \item \emph{retorna}


\end{enumerate}

\questionTitle{24}{Mecânica da execução: Atribuição vs. Avaliação}
\begin{flushright}
  \footnotesize
  \textit{\textbf{1: Paradigmas de programação}}\\
  \textit{c01-mecanismo-reducao-01}\\

\end{flushright}
No paradigma imperativo (como na linguagem C), a computação avança predominantemente através da \textbf{execução de comandos} que alteram o estado da memória. Em contraste, qual é o mecanismo nuclear que impulsiona a computação em um programa funcional puro como Haskell?

\begin{enumerate}
  \item A avaliação e redução de expressões através da aplicação de funções, substituindo parâmetros até que se atinja um valor normal (final).

  \item O despacho dinâmico de métodos orientados a eventos, desencadeados por interrupções do sistema operacional.

  \item A modificação sequencial de valores armazenados em registradores de variáveis estáticas no escopo global do módulo compilado.

  \item A validação de axiomas através de um motor de inferência lógica operando sobre um banco de dados de regras.

  \item A mutação explícita de estruturas de dados iteráveis através da alocação manual de ponteiros em tempo de execução.


\end{enumerate}

\questionTitle{25}{Inferência em discriminante de equação quadrática}
\begin{flushright}
  \footnotesize
  \textit{\textbf{7: Polimorfismo}}\\
  \textit{cap07-inferencia-possui-raizes}\\

\end{flushright}
Considere:

\begin{minted}{haskell}
possuiRaizesReais a b c = b^2 - 4*a*c >= 0
\end{minted}

Qual assinatura é a mais geral?

\begin{enumerate}
  \item \mintinline{haskell}{possuiRaizesReais :: Ord a => a -> a -> a -> Bool}
  \item \mintinline{haskell}{possuiRaizesReais :: Bool -> Bool -> Bool -> Bool}
  \item \mintinline{haskell}{possuiRaizesReais :: Floating a => a -> a -> a -> [a]}
  \item \mintinline{haskell}{possuiRaizesReais :: Num a => a -> a -> a -> Bool}
  \item \mintinline{haskell}{possuiRaizesReais :: (Num a, Ord a) => a -> a -> a -> Bool}

\end{enumerate}

\questionTitle{26}{Identificação de chamada de cauda}
\begin{flushright}
  \footnotesize
  \textit{\textbf{9: Funções recursivas}}\\
  \textit{c09-cauda-identificacao-01}\\

\end{flushright}
Qual das funções abaixo apresenta a chamada recursiva em posição de cauda?

\begin{enumerate}
  \item \begin{haskellcode}
produtoIntervalo m n
  | m > n     = 1
  | otherwise = m * produtoIntervalo (m + 1) n
\end{haskellcode}

  \item \begin{haskellcode}
pdois' x y
  | x == 0 = y
  | x > 0  = pdois' (x - 1) (2 * y)
\end{haskellcode}

  \item \begin{haskellcode}
fib n
  | n == 0 = 0
  | n == 1 = 1
  | n > 1  = fib (n - 2) + fib (n - 1)
\end{haskellcode}

  \item \begin{haskellcode}
potDois n
  | n == 0 = 1
  | n > 0  = 2 * potDois (n - 1)
\end{haskellcode}

  \item \begin{haskellcode}
fatorial n
  | n == 0 = 1
  | n > 0  = n * fatorial (n - 1)
\end{haskellcode}


\end{enumerate}

\questionTitle{27}{Instanciação de variáveis de tipo em map}
\begin{flushright}
  \footnotesize
  \textit{\textbf{7: Polimorfismo}}\\
  \textit{cap07-param-map-instanciacao}\\

\end{flushright}
Considere a assinatura:

\begin{minted}{haskell}
map :: (a -> b) -> [a] -> [b]
\end{minted}

Na expressão abaixo, quais instanciações de \haskellinline{a} e \haskellinline{b} são usadas?

\begin{minted}{haskell}
map even [1, 2, 3, 4]
\end{minted}

\begin{enumerate}
  \item \haskellinline{a} é instanciado como \haskellinline{Bool} e \haskellinline{b} como \haskellinline{Integer}, pois \haskellinline{even} retorna números pares.
  \item \haskellinline{a} é instanciado como um tipo integral, usualmente \haskellinline{Integer} no contexto padrão do GHCi, e \haskellinline{b} é instanciado como \haskellinline{Bool}.
  \item \haskellinline{a} e \haskellinline{b} são ambos instanciados como \haskellinline{Bool}, pois o resultado de \haskellinline{even} é booleano.
  \item A expressão não é tipável, pois \haskellinline{map} só aceita listas de funções.
  \item \haskellinline{a} e \haskellinline{b} permanecem totalmente independentes, pois \haskellinline{map} nunca usa o tipo da função recebida.

\end{enumerate}

\questionTitle{28}{Sintaxe explícita e regra de layout}
\begin{flushright}
  \footnotesize
  \textit{\textbf{3: Expressões e definições}}\\
  \textit{c03-layout-avaliacao-02}\\

\end{flushright}
Qual trecho é equivalente à expressão com delimitadores explícitos \haskellinline|let { a = 1; b = 2 } in a + b|?

\begin{enumerate}
  \item \begin{haskellcode}
let a = 1; b = 2; a + b
\end{haskellcode}

  \item \begin{haskellcode}
let a = 1
in b = 2
in a + b
\end{haskellcode}

  \item \begin{haskellcode}
a = 1
b = 2
in a + b
\end{haskellcode}

  \item \begin{haskellcode}
let a = 1
  b = 2
in a + b
\end{haskellcode}

  \item \begin{haskellcode}
let a = 1
    b = 2
in a + b
\end{haskellcode}


\end{enumerate}

\questionTitle{29}{Produto de intervalo vazio}
\begin{flushright}
  \footnotesize
  \textit{\textbf{9: Funções recursivas}}\\
  \textit{c09-projeto-produto-intervalo-01}\\

\end{flushright}
Na definição de \haskellinline|produtoIntervalo m n|, adotou-se a convenção de que o produto de um intervalo vazio é \haskellinline|1|. Qual caso base implementa essa convenção?

\begin{enumerate}
  \item \haskellinline|m == n = 1|

  \item \haskellinline|m < n = 1|

  \item \haskellinline|m == 0 = n|

  \item \haskellinline|m > n = 1|

  \item \haskellinline|n == 0 = 0|


\end{enumerate}

\questionTitle{30}{Soma de quantidade fixa de valores}
\begin{flushright}
  \footnotesize
  \textit{\textbf{10: Ações de E/S recursivas}}\\
  \textit{c10-rastreamento-04}\\

\end{flushright}
A ação abaixo é chamada como \haskellinline|lerESomar 3|. Depois disso, o usuário digita exatamente os valores \haskellinline|4, 5, 6|, um por linha.

\begin{haskellcode}
lerESomar :: Integer -> IO Integer
lerESomar n
  | n <= 0 = return 0
  | n > 0  = do
      x <- readLn
      s <- lerESomar (n - 1)
      return (x + s)
\end{haskellcode}

Qual valor é produzido pela ação?

\begin{enumerate}
  \item 15

  \item 14

  \item 18

  \item 3

  \item 0


\end{enumerate}

\questionTitle{31}{Valor produzido por uma ação não é impresso automaticamente}
\begin{flushright}
  \footnotesize
  \textit{\textbf{10: Ações de E/S recursivas}}\\
  \textit{c10-return-tipos-04}\\

\end{flushright}
Analise o programa:

\begin{haskellcode}
main :: IO ()
main = do
  return 42
  putStrLn "fim"
\end{haskellcode}

O que ocorre quando \haskellinline|main| é executado?

\begin{enumerate}
  \item A execução termina em \haskellinline|return 42| e a linha \haskellinline|putStrLn "fim"| nunca é executada.

  \item O número \haskellinline|42| é impresso e, em seguida, a palavra \haskellinline|"fim"| é impressa.

  \item A ação \haskellinline|return 42| é executada e seu resultado é descartado; depois, \haskellinline|"fim"| é impresso.

  \item O programa não compila, pois uma ação \haskellinline|IO Int| não pode aparecer antes de uma ação \haskellinline|IO ()| em um bloco \haskellinline|do|.

  \item O valor \haskellinline|42| é armazenado automaticamente em uma variável chamada \haskellinline|it| dentro de \haskellinline|main|.


\end{enumerate}

\questionTitle{32}{Consequência da falta de cobertura}
\begin{flushright}
  \footnotesize
  \textit{\textbf{5: Expressão condicional}}\\
  \textit{c05-guardas-exaustividade-01}\\

\end{flushright}
O que acontece se uma função com guardas é chamada com uma entrada para a qual nenhuma guarda resulta em \haskellinline|True|, e não existe outra equação aplicável?

\begin{enumerate}
  \item A função passa a retornar automaticamente \haskellinline|undefined|.

  \item A função retorna o valor da última expressão escrita na definição.

  \item O compilador sempre rejeita o programa com erro de compilação.

  \item O avaliador repete a verificação das guardas até alguma se tornar verdadeira.

  \item Ocorre um erro em tempo de execução, porque a definição não cobre aquela entrada.


\end{enumerate}

\questionTitle{33}{Fundação teórica: Cálculo Lambda}
\begin{flushright}
  \footnotesize
  \textit{\textbf{1: Paradigmas de programação}}\\
  \textit{c01-historia-evolucao-01}\\

\end{flushright}
Na década de 1930, antes mesmo da existência dos computadores eletrônicos construídos em silício, foi desenvolvido um sistema formal matemático que provou ser tão poderoso e universal quanto a Máquina de Turing. Este sistema lançou a fundação teórica de toda a programação funcional moderna.

Qual é o nome deste sistema matemático e quem foi o seu criador?

\begin{enumerate}
  \item O Sistema de Inferência de Tipos de Primeira Ordem, desenvolvido por Robin Milner.

  \item A Teoria da Avaliação Preguiçosa (\emph{Lazy Evaluation}), formalizada por David Turner.

  \item A Álgebra Booleana, desenvolvida por George Boole.

  \item O Cálculo Lambda, desenvolvido por Alonzo Church.

  \item A Arquitetura Pura, desenvolvida por John von Neumann.


\end{enumerate}

\questionTitle{34}{Comparação entre soma tardia e soma acumulada}
\begin{flushright}
  \footnotesize
  \textit{\textbf{10: Ações de E/S recursivas}}\\
  \textit{c10-acumuladores-cauda-04}\\

\end{flushright}
Compare as duas definições:

\begin{haskellcode}
-- Versão A
lerA :: IO Integer
lerA = do
  n <- readLn
  if n == 0
    then return 0
    else do
      s <- lerA
      return (n + s)

-- Versão B
lerB :: Integer -> IO Integer
lerB total = do
  n <- readLn
  if n == 0
    then return total
    else lerB (total + n)
\end{haskellcode}

Qual alternativa expressa melhor a diferença estrutural entre elas?

\begin{enumerate}
  \item As duas versões têm exatamente a mesma estrutura de chamadas, diferindo apenas nos nomes das variáveis.

  \item A versão A é pura e a versão B é impura, pois apenas B recebe parâmetro.

  \item A versão A lê todos os valores antes de testar o sentinela; a versão B testa o sentinela antes de ler.

  \item A versão A deixa uma soma pendente após a chamada recursiva; a versão B incorpora o valor lido ao acumulador antes da próxima chamada.

  \item A versão B ignora os valores lidos, pois não há variável ligada por \haskellinline|<-| no ramo recursivo.


\end{enumerate}

\questionTitle{35}{Precedência da aplicação sobre operadores}
\begin{flushright}
  \footnotesize
  \textit{\textbf{3: Expressões e definições}}\\
  \textit{c03-precedencia-associatividade-02}\\

\end{flushright}
A expressão \haskellinline|sqrt 16 + 3 * 2| ilustra a precedência da aplicação de função em relação aos operadores infixos. Qual alternativa descreve corretamente sua avaliação?

\begin{enumerate}
  \item Ela é lida como \haskellinline|sqrt (16 + 3) * 2|, resultando em aproximadamente \haskellinline|8.7178|.

  \item Ela é inválida porque operadores aritméticos não podem ser misturados com aplicação prefixa.

  \item Ela é lida como \haskellinline|sqrt (16 + 3 * 2)|, resultando em aproximadamente \haskellinline|4.6904|.

  \item Ela é lida como \haskellinline|(sqrt 16) + (3 * 2)|, resultando em \haskellinline|10.0|.

  \item Ela é inválida porque \haskellinline|sqrt| exige parênteses em volta de \haskellinline|16|.


\end{enumerate}

\questionTitle{36}{Entrada não é uma função pura}
\begin{flushright}
  \footnotesize
  \textit{\textbf{8: Programas interativos}}\\
  \textit{c08-modelo-io-04}\\

\end{flushright}
Por que uma função que lê uma linha do teclado não pode ter simplesmente o tipo \haskellinline|String|?

\begin{enumerate}
  \item Porque toda entrada digitada pelo usuário é sempre um único caractere, então o tipo correto seria \haskellinline|Char|.

  \item Porque o valor lido depende de uma interação externa. Haskell representa essa interação como uma ação do tipo \haskellinline|IO String|, não como uma \haskellinline|String| pura.

  \item Porque \haskellinline|String| é um tipo numérico e não pode ser usado para texto vindo do terminal.

  \item Porque \haskellinline|String| só pode representar literais escritos diretamente no código-fonte.

  \item Porque o compilador converte automaticamente toda leitura do teclado para \haskellinline|Int|.


\end{enumerate}

\questionTitle{37}{Turma vazia no fechamento de notas}
\begin{flushright}
  \footnotesize
  \textit{\textbf{10: Ações de E/S recursivas}}\\
  \textit{c10-bordas-robustez-04}\\

\end{flushright}
Em um programa de fechamento de notas, suponha que os dados dos alunos sejam lidos em uma lista e que o percentual de aprovados seja calculado por uma função pura que divide pelo tamanho da lista. Qual cuidado é necessário?

\begin{enumerate}
  \item A função deve chamar \haskellinline|readLn| novamente se a lista estiver vazia, mesmo sendo uma função pura.

  \item Nenhum cuidado é necessário, pois \haskellinline|length []| é automaticamente convertido para \haskellinline|1| em divisões.

  \item É necessário tratar explicitamente a turma vazia, evitando divisão por zero e definindo uma mensagem ou convenção adequada para esse caso.

  \item O problema só ocorre se as matrículas forem negativas; o tamanho da lista não interfere nos percentuais.

  \item A lista vazia deve ser substituída por uma lista infinita, para que a média da turma exista.


\end{enumerate}

\questionTitle{38}{O uso exploratório do :help}
\begin{flushright}
  \footnotesize
  \textit{\textbf{2: Ambiente de desenvolvimento}}\\
  \textit{c02-ghci-inspecao-02}\\

\end{flushright}
Qual das seguintes opções descreve a função do comando \texttt{:help} (ou \texttt{:?}) no ambiente interativo GHCi?

\begin{enumerate}
  \item Ele formata o arquivo atual para corrigir falhas de indentação antes do carregamento.

  \item Ele aciona o depurador passo a passo (\emph{debugger}) para ajudar a encontrar falhas lógicas em tempo de execução.

  \item Ele analisa o código que está sendo digitado e sugere o preenchimento automático das funções em tempo real.

  \item Ele conecta o terminal à internet para buscar soluções de erros de compilação no fórum oficial da linguagem.

  \item Ele exibe um manual interno listando todos os comandos especiais do GHCi e suas respectivas funções.


\end{enumerate}

\questionTitle{39}{Construção da lista durante o retorno da recursão}
\begin{flushright}
  \footnotesize
  \textit{\textbf{10: Ações de E/S recursivas}}\\
  \textit{c10-listas-ordem-02}\\

\end{flushright}
Analise a versão sem acumulador:

\begin{haskellcode}
lerLista :: IO [Integer]
lerLista = do
  x <- readLn
  if x == 0
    then return []
    else do
      resto <- lerLista
      return (x:resto)
\end{haskellcode}

Se o usuário digita \haskellinline|4|, depois \haskellinline|9|, e depois \haskellinline|0|, qual lista é produzida?

\begin{enumerate}
  \item \haskellinline|[0,9,4]|, pois a recursão reconstrói a lista de trás para frente incluindo o caso base.

  \item \haskellinline|[9,4]|, pois cada novo elemento é colocado no início da lista final.

  \item \haskellinline|[]|, pois \haskellinline|return []| descarta os valores anteriores.

  \item \haskellinline|[4,9,0]|, pois o sentinela também faz parte da lista retornada.

  \item \haskellinline|[4,9]|, pois cada chamada adiciona seu valor à frente da lista produzida pelo restante da recursão.


\end{enumerate}

\questionTitle{40}{A ação getContents}
\begin{flushright}
  \footnotesize
  \textit{\textbf{8: Programas interativos}}\\
  \textit{c08-entrada-saida-04}\\

\end{flushright}
Qual alternativa descreve corretamente a ação \haskellinline|getContents|?

\begin{enumerate}
  \item \haskellinline|getContents :: IO Char| lê apenas o primeiro caractere digitado pelo usuário.

  \item \haskellinline|getContents :: String| é uma string pura contendo o texto digitado antes da execução do programa.

  \item \haskellinline|getContents :: FilePath -> IO String| exige sempre o nome de um arquivo como argumento.

  \item \haskellinline|getContents :: IO String| é uma ação que lê todo o conteúdo disponível na entrada padrão como uma única string, usualmente de forma preguiçosa.

  \item \haskellinline|getContents :: IO ()| é uma ação que descarta toda a entrada padrão.


\end{enumerate}

\questionTitle{41}{A saída segura do ambiente interativo}
\begin{flushright}
  \footnotesize
  \textit{\textbf{2: Ambiente de desenvolvimento}}\\
  \textit{c02-fluxo-trabalho-02}\\

\end{flushright}
Após terminar o desenvolvimento, testar exaustivamente o arquivo \texttt{.hs} e salvar o progresso, qual comando do GHCi é o responsável por encerrar a sessão interativa de forma limpa, devolvendo o controle ao terminal do sistema operacional?

\begin{enumerate}
  \item O comando \texttt{:end} (sem abreviação suportada).

  \item O comando \texttt{:close} (ou a abreviação \texttt{:c}).

  \item O comando \texttt{:quit} (ou a abreviação \texttt{:q}).

  \item O comando \texttt{:exit} (ou a abreviação \texttt{:e}).

  \item O comando \texttt{:stop} (ou a abreviação \texttt{:s}).


\end{enumerate}

\questionTitle{42}{Tupla, lista e agrupamento}
\begin{flushright}
  \footnotesize
  \textit{\textbf{6: Estruturas de dados básicas}}\\
  \textit{c06-tuplas-sintaxe-01}\\

\end{flushright}
Qual das expressões abaixo representa uma tupla de exatamente dois elementos?

\begin{enumerate}
  \item \haskellinline{[42, "ok"]}

  \item \haskellinline{(42)}

  \item \haskellinline{(42, "ok")}

  \item \haskellinline{42, "ok"}

  \item \haskellinline{("ok")}


\end{enumerate}

\questionTitle{43}{Associatividade da aplicação de função}
\begin{flushright}
  \footnotesize
  \textit{\textbf{3: Expressões e definições}}\\
  \textit{c03-aplicacao-notacoes-02}\\

\end{flushright}
A aplicação de função em Haskell associa à esquerda. Assim, como a expressão \haskellinline|logBase 2 8| deve ser lida?

\begin{enumerate}
  \item Como \haskellinline|logBase (2 8)|, pois o argumento mais à direita é agrupado primeiro.

  \item Como uma expressão inválida, pois \haskellinline|logBase| exige que os argumentos sejam separados por vírgula.

  \item Como \haskellinline|(logBase 2 8)|, mas com \haskellinline|2| e \haskellinline|8| formando implicitamente uma tupla.

  \item Como \haskellinline|logBase (2, 8)|, pois funções de dois argumentos recebem tuplas por padrão.

  \item Como \haskellinline|(logBase 2) 8|: primeiro obtém-se a função logarítmica de base \haskellinline|2|, depois ela é aplicada a \haskellinline|8|.


\end{enumerate}

\questionTitle{44}{Composição de take, drop e sum}
\begin{flushright}
  \footnotesize
  \textit{\textbf{6: Estruturas de dados básicas}}\\
  \textit{c06-listas-operacoes-01}\\

\end{flushright}
Analise a expressão abaixo.
\begin{minted}{haskell}
let xs = [2..14]
    ys = take 5 (drop 2 xs)
in sum ys
\end{minted}
Qual é o valor final calculado?

\begin{enumerate}
  \item \haskellinline{ 20 }

  \item \haskellinline{ 104 }

  \item A expressão resulta em erro.

  \item \haskellinline{ 35 }

  \item \haskellinline{ 30 }


\end{enumerate}

\questionTitle{45}{Inferência de tipo de not True}
\begin{flushright}
  \footnotesize
  \textit{\textbf{4: Tipos de Dados}}\\
  \textit{c04-inferencia-basica-01}\\

\end{flushright}
Qual é o tipo inferido para a expressão \haskellinline{not True}?

\begin{enumerate}
  \item \haskellinline{Boolean}

  \item \haskellinline{False}

  \item \haskellinline{Bool}

  \item \haskellinline{Bool -> Bool}

  \item \haskellinline{True}


\end{enumerate}

\questionTitle{46}{Servidor de Linguagem Haskell (HLS)}
\begin{flushright}
  \footnotesize
  \textit{\textbf{2: Ambiente de desenvolvimento}}\\
  \textit{c02-ecossistema-ferramentas-02}\\

\end{flushright}
Ao configurar o editor Visual Studio Code (VS Code) para desenvolvimento Haskell, instala-se o \emph{Haskell Language Server} (HLS). Qual benefício arquitetural essa ferramenta adiciona ao editor?

\begin{enumerate}
  \item Sincroniza os arquivos locais automaticamente com as respostas exigidas no Moodle por meio de \emph{commits} contínuos.

  \item Fornece funcionalidades de \emph{IDE} inteligentes, como autocompletar, verificação de erros em tempo real e exibição de tipos ao passar o cursor sobre funções.

  \item Compila o código-fonte em segundo plano para enviar o binário final diretamente a plataformas de hospedagem em nuvem.

  \item Altera as fontes tipográficas e o esquema de cores do editor para padronizar a exibição visual de mônadas e funtores.

  \item Obriga o VS Code a rodar em um servidor web remoto, permitindo que múltiplos alunos programem no mesmo arquivo simultaneamente.


\end{enumerate}

\questionTitle{47}{Comentários aninhados e código efetivo}
\begin{flushright}
  \footnotesize
  \textit{\textbf{3: Expressões e definições}}\\
  \textit{c03-comentarios-valores-02}\\

\end{flushright}
Considere o trecho a seguir.

\begin{haskellcode}
x = 10
{-
y = 20
  {- z = 30 -}
w = 40
-}
r = x + 1
\end{haskellcode}

Quais definições permanecem como código efetivo para o compilador?

\begin{enumerate}
  \item Apenas \haskellinline|x| e \haskellinline|r|.

  \item \haskellinline|x|, \haskellinline|y|, \haskellinline|w| e \haskellinline|r|.

  \item Apenas \haskellinline|r|, pois \haskellinline|x| fica dentro do comentário de bloco.

  \item Nenhuma definição, pois comentários de bloco aninhados invalidam o arquivo.

  \item \haskellinline|x|, \haskellinline|z| e \haskellinline|r|.


\end{enumerate}

\questionTitle{48}{Validação após leitura numérica}
\begin{flushright}
  \footnotesize
  \textit{\textbf{8: Programas interativos}}\\
  \textit{c08-conversao-validacao-04}\\

\end{flushright}
Considere a ação:

\begin{minted}{haskell}
main :: IO ()
main = do
  putStrLn "Digite uma idade:"
  idade <- readLn :: IO Int
  if idade < 0
    then putStrLn "Idade inválida!"
    else print idade
\end{minted}

Qual afirmação é correta?

\begin{enumerate}
  \item O programa valida idades negativas depois de converter a entrada para \haskellinline|Int|. Entretanto, uma entrada textual que não represente um inteiro ainda pode causar exceção de conversão.

  \item O programa valida qualquer entrada inválida, incluindo textos como \haskellinline|"abc"|, sem risco de exceção.

  \item O programa não compila, pois \haskellinline|readLn| não pode receber anotação de tipo.

  \item O ramo \haskellinline|else| transforma automaticamente a idade em \haskellinline|String|.

  \item O teste \haskellinline|idade < 0| é executado antes da leitura da entrada.


\end{enumerate}

\questionTitle{49}{Operações pendentes na pilha}
\begin{flushright}
  \footnotesize
  \textit{\textbf{9: Funções recursivas}}\\
  \textit{c09-rastreamento-pilha-01}\\

\end{flushright}
Na definição
\begin{haskellcode}
fatorial n
  | n == 0 = 1
  | n > 0  = n * fatorial (n - 1)
\end{haskellcode}
por que a chamada \haskellinline|fatorial 4| deixa operações pendentes?

\begin{enumerate}
  \item Porque a multiplicação é executada antes da chamada recursiva terminar.

  \item Porque o caso base de \haskellinline|fatorial| nunca é alcançado para entradas positivas.

  \item Porque toda função com guardas é automaticamente não terminante.

  \item Porque cada chamada ainda precisa multiplicar o resultado da chamada recursiva pelo valor atual de \haskellinline|n|.

  \item Porque \haskellinline|Integer| força o uso de listas intermediárias.


\end{enumerate}

\questionTitle{50}{Paradigmas - Princípio Nuclear (Legado Refatorado)}
\begin{flushright}
  \footnotesize
  \textit{\textbf{1: Paradigmas de programação}}\\
  \textit{c01-conceitos-paradigmas-02}\\

\end{flushright}
Ao contrastarmos os paradigmas dominantes, o \textbf{princípio nuclear} que separa a programação imperativa da programação declarativa (funcional) reside no foco da estruturação do código. Qual é essa diferença?

\begin{enumerate}
  \item A programação imperativa exige a declaração estrita de tipos em tempo de compilação, enquanto a declarativa permite que os tipos sejam inferidos dinamicamente durante a execução.

  \item A programação imperativa foi projetada exclusivamente para arquiteturas mononucleares (single-core), enquanto a declarativa foi projetada nativamente para paralelismo em GPUs.

  \item A programação imperativa abstrai o gerenciamento de memória através de \emph{Garbage Collectors}, enquanto a declarativa exige alocação explícita com \cinline|malloc|.

  \item A programação imperativa foca em descrever \emph{como} o computador deve executar os passos (controle de fluxo e estado), enquanto a declarativa foca em descrever \emph{o que} é o resultado desejado (lógica e propriedades).

  \item A programação imperativa é estruturada em torno de funções matemáticas puras, enquanto a declarativa é estruturada em torno de objetos que trocam mensagens.


\end{enumerate}

\questionTitle{51}{Eficiência e recomputação}
\begin{flushright}
  \footnotesize
  \textit{\textbf{9: Funções recursivas}}\\
  \textit{c09-eficiencia-fibonacci-01}\\

\end{flushright}
Ao comparar a definição ingênua de Fibonacci com uma versão usando acumuladores, qual é a principal diferença de eficiência discutida no capítulo?

\begin{enumerate}
  \item A versão com acumuladores deixa de usar casos base.

  \item As duas versões sempre têm exatamente o mesmo número de chamadas.

  \item A versão com acumuladores só funciona para \haskellinline|fib 0| e \haskellinline|fib 1|.

  \item A versão ingênua é mais rápida porque usa duas chamadas recursivas.

  \item A versão ingênua recalcula subproblemas repetidamente, enquanto a versão com acumuladores carrega resultados parciais de forma mais direta.


\end{enumerate}

\questionTitle{52}{Anatomia de uma definição em arquivo}
\begin{flushright}
  \footnotesize
  \textit{\textbf{3: Expressões e definições}}\\
  \textit{c03-arquivos-identificadores-02}\\

\end{flushright}
Considere a definição em um arquivo fonte:

\begin{haskellcode}
dobro n = 2 * n
\end{haskellcode}

Qual alternativa identifica corretamente seus componentes?

\begin{enumerate}
  \item A definição é inválida porque funções de um argumento exigem vírgula depois do nome.

  \item \haskellinline|dobro n| é o nome completo da função, e por isso ela não recebe argumentos.

  \item \haskellinline|dobro| é o identificador definido, \haskellinline|n| é o parâmetro e \haskellinline|2 * n| é o corpo da definição.

  \item \haskellinline|2 * n| é uma assinatura de tipo.

  \item \haskellinline|n| é uma definição global independente criada automaticamente.


\end{enumerate}

\questionTitle{53}{Tipo de um bloco do}
\begin{flushright}
  \footnotesize
  \textit{\textbf{8: Programas interativos}}\\
  \textit{c08-tipos-acoes-04}\\

\end{flushright}
Qual é o tipo da ação abaixo?

\begin{minted}{haskell}
acao = do
  putStrLn "Digite uma palavra:"
  getLine
\end{minted}

\begin{enumerate}
  \item \haskellinline|IO String|, pois a última ação do bloco é \haskellinline|getLine|.

  \item \haskellinline|String|, pois o bloco \haskellinline|do| extrai automaticamente o valor produzido por \haskellinline|getLine|.

  \item \haskellinline|String -> IO String|, pois a primeira linha imprime uma mensagem.

  \item \haskellinline|IO ()|, pois o bloco contém uma chamada a \haskellinline|putStrLn|.

  \item \haskellinline|()|, pois o valor lido por \haskellinline|getLine| é descartado.


\end{enumerate}

\questionTitle{54}{Benefício de testar funções puras separadamente}
\begin{flushright}
  \footnotesize
  \textit{\textbf{10: Ações de E/S recursivas}}\\
  \textit{c10-separacao-puro-02}\\

\end{flushright}
Ao projetar um programa interativo em Haskell, por que é desejável separar a coleta de dados em \haskellinline|IO| do processamento puro?

\begin{enumerate}
  \item Porque funções puras imprimem seus resultados automaticamente no GHCi e no programa compilado.

  \item Porque ações de E/S não podem ter assinaturas de tipo explícitas.

  \item Porque o compilador exige que todo bloco \haskellinline|do| contenha exatamente uma função pura.

  \item Porque funções puras podem ser testadas e reutilizadas sem depender de entrada e saída reais.

  \item Porque funções puras são as únicas que podem chamar \haskellinline|readLn| com segurança.


\end{enumerate}

\questionTitle{55}{Quando guardas são mais idiomáticas}
\begin{flushright}
  \footnotesize
  \textit{\textbf{5: Expressão condicional}}\\
  \textit{c05-if-vs-guardas-01}\\

\end{flushright}
Em qual situação o uso de guardas tende a ser mais idiomático e legível do que \haskellinline|if-then-else| aninhado?

\begin{enumerate}
  \item Quando se deseja evitar qualquer comparação booleana.

  \item Quando os ramos do \haskellinline|if| têm tipos diferentes.

  \item Quando a função trata vários casos sequenciais, cada um coberto por uma guarda.

  \item Quando há apenas um teste booleano simples, com dois resultados possíveis.

  \item Quando a função não recebe argumentos.


\end{enumerate}

\questionTitle{56}{Aplicações de missão crítica (Legado Refatorado)}
\begin{flushright}
  \footnotesize
  \textit{\textbf{1: Paradigmas de programação}}\\
  \textit{c01-engenharia-software-01}\\

\end{flushright}
Você foi encarregado de arquitetar o \emph{backend} de um sistema bancário de larga escala. O sistema processa milhões de transações diárias que passam por validações, conversões e cálculos de risco. A \textbf{correção matemática} e a \textbf{auditabilidade} do fluxo de dados são características críticas inegociáveis.

Entre as opções, qual paradigma de programação fornece as abstrações arquiteturais mais seguras para este domínio específico e por quê?

\begin{enumerate}
  \item O paradigma puramente orientado a objetos, porque ocultar transações dentro de objetos de estado opaco garante que auditores externos não consigam corromper os dados acidentalmente.

  \item O paradigma lógico, porque todo o ecossistema bancário pode ser facilmente abstraído como um banco de dados Prolog e resolvido instantaneamente via inferência de axiomas.

  \item O paradigma funcional, porque a imutabilidade estrita e a ausência de efeitos colaterais tornam o fluxo de dados totalmente explícito e previsível, facilitando a auditoria e as provas lógicas.

  \item A escolha do paradigma é completamente irrelevante para características de software como auditabilidade e correção matemática; essas métricas dependem exclusivamente da linguagem SQL utilizada.

  \item O paradigma imperativo procedural, porque a manipulação direta e mutável do estado da memória é a única forma técnica de garantir a performance de I/O exigida por bancos de dados modernos.


\end{enumerate}

\questionTitle{57}{Avaliação com where e dependência local}
\begin{flushright}
  \footnotesize
  \textit{\textbf{3: Expressões e definições}}\\
  \textit{c03-definicoes-where-02}\\

\end{flushright}
Considere a definição:

\begin{haskellcode}
f x = a + b
  where
    a = x + 1
    b = a * 2
\end{haskellcode}

Qual é o valor de \haskellinline|f 3|?

\begin{enumerate}
  \item \haskellinline|4|.

  \item \haskellinline|8|.

  \item \haskellinline|12|.

  \item Ocorre erro, pois \haskellinline|b| não pode usar \haskellinline|a| definido no mesmo bloco \haskellinline|where|.

  \item \haskellinline|7|.


\end{enumerate}

\questionTitle{58}{Definição formal de transparência referencial}
\begin{flushright}
  \footnotesize
  \textit{\textbf{1: Paradigmas de programação}}\\
  \textit{c01-transparencia-referencial-01}\\

\end{flushright}
Na teoria de linguagens de programação funcionais, o que significa afirmar que uma dada expressão possui \textbf{transparência referencial}?

\begin{enumerate}
  \item Significa que a expressão pode ser substituída diretamente pelo seu valor computado correspondente em qualquer ponto do programa, sem alterar o comportamento global do sistema.

  \item Significa que a função tem permissão explícita para ler variáveis do escopo global mutável, desde que não as modifique durante a execução.

  \item Significa que o sistema de inferência de tipos consegue deduzir a assinatura da função mesmo que o programador omita essa informação no cabeçalho.

  \item Significa que o código-fonte da função não pode ser ofuscado pelo compilador, permanecendo \emph{transparente} e legível para fins de auditoria no binário.

  \item Significa que a função permite que referências de ponteiros sejam passadas como argumentos em vez de cópias de valores, economizando alocação de memória.


\end{enumerate}

\questionTitle{59}{Leitura de condicional aninhada}
\begin{flushright}
  \footnotesize
  \textit{\textbf{5: Expressão condicional}}\\
  \textit{c05-if-aninhamento-01}\\

\end{flushright}
Considere a definição:
\begin{haskellcode}
sinal n = if n < 0 then -1 else if n > 0 then 1 else 0
\end{haskellcode}
Qual é o resultado de \haskellinline|sinal 0|?

\begin{enumerate}
  \item \haskellinline|-1|

  \item Ocorre um erro, pois o segundo \haskellinline|if| não tem \haskellinline|else|.

  \item \haskellinline|False|

  \item \haskellinline|1|

  \item \haskellinline|0|


\end{enumerate}

\questionTitle{60}{Domínio de uma função recursiva}
\begin{flushright}
  \footnotesize
  \textit{\textbf{9: Funções recursivas}}\\
  \textit{c09-fundamentos-dominio-01}\\

\end{flushright}
Ao definir uma função recursiva sobre números naturais em Haskell, por que é importante explicitar o domínio ou tratar entradas fora dele?

\begin{enumerate}
  \item Porque funções recursivas só podem ser usadas com o tipo \haskellinline|Natural|, nunca com \haskellinline|Integer|.

  \item Porque todo caso recursivo precisa ter exatamente dois casos base.

  \item Porque a função pode receber valores fora do domínio matemático pretendido, e isso pode causar falhas de cobertura ou chamadas que nunca alcançam o caso base.

  \item Porque o compilador sempre rejeita funções recursivas que não tratam números negativos.

  \item Porque Haskell não permite definir funções recursivas sem testar todos os inteiros possíveis explicitamente.


\end{enumerate}

\questionTitle{61}{Visibilidade em guardas e resultados}
\begin{flushright}
  \footnotesize
  \textit{\textbf{5: Expressão condicional}}\\
  \textit{c05-where-escopo-01}\\

\end{flushright}
Em uma equação com guardas e uma cláusula \haskellinline|where|, onde são visíveis as definições locais?

\begin{enumerate}
  \item Em toda a função, inclusive em outras equações com o mesmo nome.

  \item Apenas nas guardas, nunca nas expressões de resultado.

  \item Apenas na primeira guarda da equação.

  \item Apenas nas expressões de resultado, nunca nas guardas.

  \item Nas guardas e nas expressões de resultado da mesma equação.


\end{enumerate}

\questionTitle{62}{Forma fundamental de uma lista}
\begin{flushright}
  \footnotesize
  \textit{\textbf{6: Estruturas de dados básicas}}\\
  \textit{c06-listas-estrutura-01}\\

\end{flushright}
Qual é a forma fundamental da lista \haskellinline{['a', 'b', 'c']} usando apenas o construtor \haskellinline{(:)} e a lista vazia?

\begin{enumerate}
  \item \haskellinline{'a' : 'b' : 'c'}

  \item \haskellinline{['a'] : ['b'] : ['c']}

  \item \haskellinline{('a', 'b', 'c')}

  \item \haskellinline{'a' : ('b' : ('c' : []))}

  \item \haskellinline{[] : 'a' : 'b' : 'c'}


\end{enumerate}

\questionTitle{63}{Cláusula else obrigatória}
\begin{flushright}
  \footnotesize
  \textit{\textbf{5: Expressão condicional}}\\
  \textit{c05-if-regras-01}\\

\end{flushright}
Em Haskell, por que a expressão abaixo é incorreta?
\begin{haskellcode}
if x > 0 then x
\end{haskellcode}

\begin{enumerate}
  \item Porque a condição \haskellinline|x > 0| deveria estar entre parênteses.

  \item Porque expressões condicionais só podem ser usadas dentro do corpo de funções nomeadas.

  \item Porque o ramo \haskellinline|then| só pode conter valores booleanos.

  \item Porque o ramo \haskellinline|then| deveria conter uma expressão mais completa do que \haskellinline|x|.

  \item Porque falta o ramo \haskellinline|else|, obrigatório para que a expressão sempre produza um valor.


\end{enumerate}

\questionTitle{64}{Valor opcional}
\begin{flushright}
  \footnotesize
  \textit{\textbf{6: Estruturas de dados básicas}}\\
  \textit{c06-maybe-conceito-01}\\

\end{flushright}
Qual opção representa corretamente a ideia central da estrutura \haskellinline{Maybe}?

\begin{enumerate}
  \item Ela representa tuplas com campos mutáveis.

  \item Ela representa explicitamente a presença ou a ausência de um valor.

  \item Ela representa listas de tamanho variável.

  \item Ela é usada apenas para números opcionais.

  \item Ela substitui o tipo \haskellinline{Bool} em expressões condicionais.


\end{enumerate}

\questionTitle{65}{Avaliação básica de guardas}
\begin{flushright}
  \footnotesize
  \textit{\textbf{5: Expressão condicional}}\\
  \textit{c05-guardas-avaliacao-01}\\

\end{flushright}
Considere a função:
\begin{haskellcode}
f x
  | x > 10 = 'A'
  | x > 5  = 'B'
  | otherwise = 'C'
\end{haskellcode}
Qual é o valor de \haskellinline|f 3|?

\begin{enumerate}
  \item \haskellinline|'C'|

  \item Ocorre um erro de guardas não exaustivas.

  \item \haskellinline|'A'|

  \item \haskellinline|'B'|

  \item \haskellinline|"C"|


\end{enumerate}

\questionTitle{66}{Classificação da forma de recursão}
\begin{flushright}
  \footnotesize
  \textit{\textbf{9: Funções recursivas}}\\
  \textit{c09-formas-classificacao-01}\\

\end{flushright}
Qual alternativa classifica corretamente as formas de recursão estudadas no capítulo?

\begin{enumerate}
  \item \haskellinline|fib| não é recursiva, pois usa soma no caso recursivo.

  \item \haskellinline|fatorial| é recursão múltipla; \haskellinline|fib| é recursão de cauda; \haskellinline|par| e \haskellinline|impar| não são recursivas.

  \item Toda função com mais de uma guarda é, necessariamente, uma função mutuamente recursiva.

  \item \haskellinline|fatorial| é recursão simples; \haskellinline|fib| é recursão múltipla; \haskellinline|par| e \haskellinline|impar| formam recursividade mútua.

  \item \haskellinline|fatorial| e \haskellinline|fib| são recursividade mútua, pois ambas têm casos base.


\end{enumerate}

\questionTitle{67}{Tipagem estática vs. dinâmica}
\begin{flushright}
  \footnotesize
  \textit{\textbf{4: Tipos de Dados}}\\
  \textit{c04-checagem-estatica-dinamica-01}\\

\end{flushright}
Qual é a principal diferença entre um sistema de tipos estático (como Haskell) e um dinâmico (como Python)?

\begin{enumerate}
  \item Não há diferença fundamental; são apenas termos diferentes para o mesmo conceito.

  \item Apenas linguagens compiladas podem ter tipagem estática.

  \item A tipagem estática é mais lenta, mas mais flexível que a dinâmica.

  \item Na tipagem estática, a verificação de tipos ocorre em tempo de compilação; na dinâmica, em tempo de execução.

  \item Na tipagem estática, os tipos podem mudar durante a execução; na dinâmica, são fixos.


\end{enumerate}

\questionTitle{68}{Ações que exibem valores retornam IO unit}
\begin{flushright}
  \footnotesize
  \textit{\textbf{10: Ações de E/S recursivas}}\\
  \textit{c10-lacos-io-04}\\

\end{flushright}
Considere a ação:

\begin{haskellcode}
mostraLista :: Show a => [a] -> IO ()
mostraLista lista
  | null lista = return ()
  | otherwise = do
      print (head lista)
      mostraLista (tail lista)
\end{haskellcode}

Por que o tipo de retorno é \haskellinline|IO ()|?

\begin{enumerate}
  \item Porque \haskellinline|return ()| interrompe o programa e impede que qualquer valor seja produzido.

  \item Porque toda ação recursiva em Haskell deve obrigatoriamente ter tipo \haskellinline|IO ()|.

  \item Porque \haskellinline|print| apaga os elementos da lista e sempre retorna a lista vazia.

  \item Porque a lista é convertida em uma tupla vazia antes de ser percorrida.

  \item Porque a ação realiza efeitos de saída, mas não produz um valor de domínio relevante ao final.


\end{enumerate}

\questionTitle{69}{Por que Int não satisfaz Fractional}
\begin{flushright}
  \footnotesize
  \textit{\textbf{7: Polimorfismo}}\\
  \textit{cap07-hierarquia-int-fracao}\\

\end{flushright}
Por que a expressão abaixo não é aceita sem conversão explícita?

\begin{minted}{haskell}
n :: Int
n = 10

n / 2
\end{minted}

\begin{enumerate}
  \item Porque literais inteiros nunca podem aparecer em expressões com divisão.
  \item Porque \haskellinline{(/)} exige \haskellinline{Fractional a => a -> a -> a}, e \haskellinline{Int} não é instância de \haskellinline{Fractional}.
  \item Porque Haskell converte automaticamente \haskellinline{Int} para \haskellinline{Double}, mas apenas fora do GHCi.
  \item Porque \haskellinline{Int} não é instância de \haskellinline{Num}.
  \item Porque \haskellinline{(/)} só existe para \haskellinline{Integer}.

\end{enumerate}


\end{multicols}
\makeanswersheet{UFOP}{BCC222: Programação Funcional}{José Romildo}{2026/1}{Primeira Prova}{09}{69}{25}{7}
\gabarito{09}{D,B,E,B,A,A,B,C,D,A,B,B,E,C,B,D,D,E,D,B,B,D,B,A,E,B,B,E,D,A,C,E,D,D,D,B,C,E,E,D,C,C,E,E,C,B,A,A,D,D,E,C,A,D,C,C,C,A,E,C,E,D,E,B,A,D,D,E,B}
\end{document}
